Now let @@M0@@. The @@M1@@-homomorphism @@M2@@ is injective by 3.1. Consequently, to prove relations V 2.4.2, we can work in @@M3@@. Since in @@M4@@ the elements @@M5@@ and @@M6@@ decompose into sums of classes of invertible sheaves, it suffices, by additivity, to prove these same relations for elements @@M7@@, classes of invertible sheaves; they can also be written:
@@M8@@
by putting them in the form V 2.4.1. They are then precisely relations V 2.3.1 and 2.3.3 defining the @@M9@@-ring structure of @@M10@@.

Let us state the following more general result:

\sgaStableTarget{sga6.e06.thm.3-3}\textbf{Theorem 3.3.} \uline{Let} @@M11@@ \uline{be a commutatively ringed topos}; \uline{write} @@M12@@ \uline{for the Grothendieck group of} @@M13@@-\uline{Modules locally free} of bounded rank on @@M14@@. \uline{Then}, @@M15@@, \uline{endowed with its canonical pre}-@@M16@@-\uline{structure}~(V 2.2 b)), \uline{is} \uline{a} @@M17@@-\uline{ring}.

To show that @@M18@@ is a @@M19@@-ring, we must prove relations V 2.4.1, and, by additivity, it suffices to prove them for elements @@M20@@ which are classes of locally free @@M21@@-Modules of bounded rank. Since the ranks of @@M22@@ and @@M23@@ are bounded, we can decompose @@M24@@ into a finite ``disjoint sum'' of opens over which @@M25@@ and @@M26@@ have constant rank. Using the commutation of @@M27@@ with finite direct sums, and putting relations V 2.4.1 in the form V 2.4.2, we see that it suffices to prove relations V 2.4.2 in each of the @@M28@@ of the opens considered, that is, it suffices to prove relations V 2.4.2 for elements @@M29@@, classes of @@M30@@-Modules @@M31@@ locally free of constant rank.

